\section{Discussion}\label{sec:discussion}

\subsection{Lessons for developers}\label{sec:disc-lessons}

The signature analogy of Section~\ref{sec:definition} doubles as a
checklist. \emph{Sign everything you act on}: every recipient, amount, fee
and refund address must be a public input or hashed into one, and the
consumer must recompute that hash (P1). \emph{Separate domains}: the
statement should commit to the chain, the consumer contract and the policy,
and roots and prior states must be checked against the consumer's own
storage (P2). \emph{Derive nonces from the statement}, reduced into the field
before they are stored (P3). \emph{Pin the key}: fix the verification key or
program ID and check the verifier's result (P4). \emph{Test with a malicious
caller}: an honest proof submitted by a second account with changed
arguments must fail; this single test, the invariant of
Section~\ref{sec:design-explore}, would have caught \DEMO{every flaw in B1
that our fuzzing harness could express}.
\TODO{D4}{back the checklist with the fix-pattern counts from Section 4}

\subsection{Toward binding by construction}\label{sec:disc-construction}

Detection is a stopgap. Binding flaws exist because the public-input vector
is an untyped array whose meaning is restated by hand in the circuit, the
verifier wrapper and the consumer. Three changes would remove most of the
class by construction. \emph{Typed statements}: circuit languages could
export a statement schema (names, types, widths, order) from which verifier
generators emit a typed consumer interface; this removes P5 and makes P1
visible. \emph{Effect commitments}: a standard effect hash over recipient,
amount, chain, contract and nonce, computed in the circuit and recomputed in
the consumer, would play the role EIP-712 plays for signatures and address
P1--P3 together. \emph{Pinned verification}: libraries could make the key or
program ID an immutable constructor argument (P4).
\DEMO{Applied to the 87 core-class findings, typed statements would have
prevented 41\% and effect commitments 63\%.}
\TODO{D6}{estimate per finding which construction would have prevented it}

\subsection{Threats to validity}\label{sec:disc-threats}

\para{Selection}
We study public audits of projects that chose to publish them. Projects
without audits may have more binding flaws, and audit firms differ in which
findings they publish. Some sources could not be fully enumerated.

\para{Labels}
Labels were proposed by an LLM and adjudicated by two collaborating raters
who discussed every item, so they are not independent annotations. We mitigate this with a
written guide fixed before labeling, a code-level second pass for every
evaluated item, \DEMO{and an independent blind relabel of 30 items
($\kappa=0.78$)}. Pattern boundaries (especially P2 versus P1 for stored
values) involve judgment, so we release all labels with their notes.

\para{Ground-truth locations}
A flaw can often be fixed at more than one place. Strict scoring requires the
location the auditors or fix authors chose, which undercounts warnings at an
equally valid place. We therefore also report mechanism-level hits.

\para{LLM judge and variance}
An LLM judge decides whether a warning describes the same mechanism. In a
re-judging test about one verdict in twenty changed. Detection itself varies
by $\pm$1--2 items between runs, which is why we report repeated runs.
% src: mvp/tool/FREEZE-v2.md known weaknesses 1, 9

\para{Leakage}
The people who built the benchmark also saw the test results of an earlier
tool version, and the B1 misses informed the final version; we therefore
count B1 as development data for it and give the full timeline in
Appendix~\ref{app:timeline}. To keep the tool independent, every tool version was built by a
separate process that could read only development data, under enforced access
rules, and the test benchmark lives in a separate repository. The round-2
labels existed before the freeze and were read by the coordinating author, but
never by the tool-building process. Every exposure is logged in the
artifact.

\subsection{Limitations}\label{sec:disc-limits}

\tool{} needs an LLM and costs one to two dollars per scan. Its tools are
syntactic, so library internals that are not vendored are blind spots, and
its knowledge base covers ten proof systems and domains; others fall back to
the generic procedure. The exploration loop targets EVM consumers: flaws
whose effect lies in a zkVM host or off chain stop at E1, and real-proof
tests were possible on only 2 of 11 development targets, because most
projects ship no proving artifacts. Most warnings therefore reach humans as
claims. Finally, binding is relative to the intended policy $\Pi$, which
\tool{} must infer from code and documentation, and it can be wrong.
